% Shared mathematical text for the monograph and standalone research note.
% MathDevMCP provenance and reversible integration mapping are in
% docs/plans/artifacts/ledh-covariance-monograph-integration-20261008-01/.
\newcommand{\bfCovOne}{\mathbf 1}
\newcommand{\bfCovNormal}{\mathcal N}
\newcommand{\bfCovDiff}{\mathrm d}
\newcommand{\bfCovSPD}{\mathbb S_{++}}
\newcommand{\bfCovLSE}{\operatorname{LSE}}
\newcommand{\bfCovSoftmax}{\operatorname{softmax}}

\section{The variance puzzle and the filtering problem}

Consider the scalar state equation
\(X_t=\phi X_{t-1}+\eta_t\), with independent
\(\eta_t\sim\bfCovNormal(0,Q)\), \(Q>0\), and \(|\phi|<1\).
At stationarity, the state variance is \(P_\infty=Q/(1-\phi^2)\).
For \(\phi=0.9\) and \(Q=1\), this is \(100/19\), much larger than one.
Nevertheless, conditional on a particular previous state \(a\), the next
state has variance exactly one. The two statements describe different
conditioning events:
\begin{equation}
\label{eq:bf-covariance-proposal-ar1-total}
 \frac{Q}{1-\phi^2}
 =\phi^2\frac{Q}{1-\phi^2}+Q.
\end{equation}
The first term on the right is uncertainty between ancestors; the second is
uncertainty within one ancestor's transition. Replacing the second by one
does not remove the first. Replacing the entire left-hand side by one would.

The same distinction matters when a Gaussian flow moves particles. A common
map centred on the whole predictive mean may correctly use the whole
predictive covariance. Separate maps centred on individual ancestors need
covariances appropriate to those conditional distributions. An importance
correction can remain mathematically valid even when the chosen maps are
poor proposals; validity of the density ratio does not guarantee usable
finite-sample error.

We work with a state \(X_t\in\R^d\), an observation \(Y_t\in\R^m\), and a parameter
\(\theta\in\R^p\). The model supplies an initial law, transition density
\(f_\theta(x\mid a)\), and observation density \(g_\theta(y\mid x)\).
For the main construction these are densities on Euclidean spaces with
respect to their stated base measures. A differentiable transition simulator
has the form \(x=F_\theta(a,\xi)\), where the law of \(\xi\) does not depend
on \(\theta\). A fixed draw of all base noises defines a deterministic
finite filtering computation.

At a particular observation, suppose the previous filtering law is
approximated by equally weighted particles \(a_1,\ldots,a_N\). Equality of
weights follows from the preceding reset, not from a claim that the previous
approximation is exact. Suppress \(t\) temporarily and define
\begin{equation}
\label{eq:bf-covariance-proposal-finite-target}
 f_j(x)=f_\theta(x\mid a_j),\qquad
 p_N^-(x)=\frac1N\sum_{j=1}^N f_j(x),\qquad
 Z_N=\int g_\theta(y\mid x)p_N^-(x)\,\bfCovDiff x .
\end{equation}
Here \(0<Z_N<\infty\) is assumed. The one-step posterior conditional on this
finite incoming cloud is
\(\pi_N^+(x)=g_\theta(y\mid x)p_N^-(x)/Z_N\).
It is distinct from the exact model posterior, because the incoming cloud
itself is an approximation. Keeping this distinction visible prevents a
one-step integration identity from becoming an unjustified claim about the
full likelihood.

For non-Euclidean or restricted state spaces, density transformations must
include any coordinate-chart Jacobian. A proposal may put mass outside the
target support only if the extended target density there is defined to be
zero. The reset must still produce states on which the next transition is
defined. Unrecorded clipping changes the computation and is not part of the
algorithm below.

\section{From conditional uncertainty to a UKF proposal}

\subsection{The covariance of the predictive mixture}

Assume the conditional first two moments are finite and available:
\(m_j=\E[X_t\mid X_{t-1}=a_j]\) and
\(C_j=\Cov(X_t\mid X_{t-1}=a_j)\).
For additive Gaussian transitions, \(m_j=F_\theta(a_j)\) and
\(C_j=Q_\theta(a_j)\). For non-additive transitions, conditional quadrature
may approximate these moments, in which case the following identity applies
to the actual mixture only when the supplied moments are exact.

\begin{proposition}[Total predictive covariance]
The mean and covariance of \(p_N^-\) are
\begin{equation}
\label{eq:bf-covariance-proposal-mixture-moments}
 m^-=\frac1N\sum_jm_j,\qquad
 P^-=\frac1N\sum_j\left[
 C_j+(m_j-m^-)(m_j-m^-)^\top\right].
\end{equation}
\end{proposition}
\begin{proof}
Let \(J\) be uniform on \(\{1,\ldots,N\}\), and condition \(X\) on
\(J=j\) using density \(f_j\). Since
\(\E[X-m_j\mid J=j]=0\), expanding the outer product of
\(X-m^-=(X-m_j)+(m_j-m^-)\) gives
\[
 \E[(X-m^-)(X-m^-)^\top\mid J=j]
 =C_j+(m_j-m^-)(m_j-m^-)^\top .
\]
The mixed terms have zero conditional expectation. Averaging over \(J\)
proves the covariance expression; the mean follows by the same conditioning.
\end{proof}

This identity reconstructs the whole predictive covariance from the actual
incoming particles. It avoids treating a separately propagated Gaussian
covariance as an exact description of their cloud. In the AR(1) example,
the between-ancestor contribution approaches \(81/19\) at stationarity,
while each conditional contribution remains one.

\subsection{The UKF observation update}

The UKF approximates observation moments under a Gaussian guide with mean
\(m^-\) and covariance \(P^-\). We require \(P^-\in\bfCovSPD^d\).
For clarity, first consider an additive observation
\(Y=h_\theta(X)+\epsilon\), with independent centred noise of covariance
\(R\in\bfCovSPD^m\). A sigma-point rule consists of standardized points
\(\zeta_r\), mean weights \(u_r\), and covariance weights \(v_r\).
Its prescribed normalization reproduces constant, linear, and quadratic
state moments. For example, a fully specified scaled unscented rule uses
frozen \(\alpha_U>0,\kappa_U,\beta_U\), with
\(\lambda_U=\alpha_U^2(d+\kappa_U)-d\) and \(d+\lambda_U>0\):
\[
 \zeta_0=0,\quad
 \zeta_k=\sqrt{d+\lambda_U}\,e_k,\quad
 \zeta_{d+k}=-\sqrt{d+\lambda_U}\,e_k
 \quad (k=1,\ldots,d),
\]
\[
 u_0=\frac{\lambda_U}{d+\lambda_U},\quad
 v_0=u_0+1-\alpha_U^2+\beta_U,\qquad
 u_r=v_r=\frac{1}{2(d+\lambda_U)}\quad(r=1,\ldots,2d).
\]
Here \(e_k\) is a coordinate basis vector. Symmetry gives
\(\sum_ru_r\zeta_r=0\); the two points on each axis give
\(\sum_rv_r\zeta_r\zeta_r^\top=I_d\).
No particular values of these tuning constants are promoted.
If \(P^-=L_-L_-^\top\) is the Cholesky factorization, its state
points are \(\chi_r=m^-+L_-\zeta_r\). Define
\begin{align}
\label{eq:bf-covariance-proposal-ukf-observation}
 \bar y_{\rm obs}&=\sum_r u_r h_\theta(\chi_r),&
 S_{\rm obs}&=\sum_r v_r(h_\theta(\chi_r)-\bar y_{\rm obs})
                  (h_\theta(\chi_r)-\bar y_{\rm obs})^\top+R.
\end{align}
Assume the realized \(S_{\rm obs}\) is positive definite and hence
invertible; the cross-covariance and gain are
\begin{align}
\label{eq:bf-covariance-proposal-ukf-cross}
 C_{xy}&=\sum_r v_r(\chi_r-m^-)
                  (h_\theta(\chi_r)-\bar y_{\rm obs})^\top,&
 K_U&=C_{xy}S_{\rm obs}^{-1}.
\end{align}
The inverse in \eqref{eq:bf-covariance-proposal-ukf-cross} requires \(S_{\rm obs}\in\bfCovSPD^m\).
Negative sigma-point covariance weights can complicate this condition, so it
must be checked for the chosen rule. The UKF guide is
\begin{equation}
\label{eq:bf-covariance-proposal-ukf-update}
 m_U^+=m^-+K_U(y-\bar y_{\rm obs}),\qquad
 P_U^+=P^- - K_US_{\rm obs}K_U^\top.
\end{equation}
We also require \(P_U^+\in\bfCovSPD^d\). Non-additive observation noise requires
an augmented rule or direct joint-moment calculation; it must not be
replaced silently by the additive formulas.

These are Gaussian-guide moments, not a theorem about the true nonlinear
posterior. Even high-order quadrature accuracy for a specified input law
does not establish that the actual posterior is Gaussian. If the supplied
conditional moments or the UKF moments are approximate, their role remains
proposal construction; the weight uses the actual model densities.

\subsection{The exact place where the UKF variance is used}

Factor \(P_U^+=L_+L_+^\top\), with positive Cholesky diagonals, and define
\begin{equation}
\label{eq:bf-covariance-proposal-global-map}
 A_G=L_+L_-^{-1},\qquad c_G=m_U^+-A_Gm^-,
 \qquad T_G(x)=A_Gx+c_G .
\end{equation}
All inverses in this note are implemented by solves. Since \(P^-\succ0\)
and \(P_U^+\succ0\), both Cholesky factors are invertible. Since both factors
are nonsingular, \(T_G\) is an invertible affine map.

\begin{proposition}[Moments of the common proposal]
If \(X\) has moments \((m^-,P^-)\), the random vector \(T_G(X)\) has
moments \((m_U^+,P_U^+)\), regardless of whether \(X\) is Gaussian.
\end{proposition}
\begin{proof}
The mean follows by substituting \(\E X=m^-\). For the covariance,
\begin{equation}
\label{eq:bf-covariance-proposal-global-covariance}
 \Cov(T_G(X))
 =L_+L_-^{-1}(L_-L_-^\top)L_-^{-\top}L_+^\top
 =L_+L_+^\top=P_U^+ .
\end{equation}
\end{proof}

Thus \(P^-\) standardizes the whole predictive spread and \(P_U^+\)
sets the proposed spread. This is an active use of UKF covariance.
It is different from inserting \(P^-\) into every ancestor-centred flow.

In the ideal linear-Gaussian case with a Gaussian predictive distribution,
the common map pushes that distribution to the exact Gaussian posterior.
With \(P^-=100/19\), observation variance \(R=1\), observation matrix one,
and predictive mean zero, the posterior variance is \(100/119\).
The common map is
\[
 T_G(x)=\frac{100}{119}y+\sqrt{\frac{19}{119}}\,x.
\]
It contracts both within-ancestor and between-ancestor spread coherently
around the whole-cloud centre. A finite Gaussian mixture need not itself
be Gaussian, so this limiting case is not a finite-\(N\) exactness claim.

After the importance update below, the covariance target changes to the
weighted particle covariance. Figure~\ref{fig:bf-covariance-proposal-covariance-roles} locates
the three operations without conflating their targets.
\begin{figure}[htbp]
\centering
\begin{tikzpicture}[>=Latex,
 box/.style={draw,rounded corners,align=center,text width=39mm,
 minimum height=17mm,font=\small},node distance=8mm]
\node[box] (guide) {Predictive moments\\\(m^-,P^-\)\\UKF guide \(m_U^+,P_U^+\)};
\node[box,right=of guide] (weight) {Draw proposals\\Correct with actual\\\(g\,p_N^-/q\)};
\node[box,right=of weight] (reset) {Weighted moments\\\(\mu_w,P_w\)\\Preserve during reset};
\draw[->] (guide)--(weight);
\draw[->] (weight)--(reset);
\end{tikzpicture}
\caption{The UKF guides sampling. Importance weights determine the moments
that resampling preserves. No arrow forces the weighted cloud back to
the UKF covariance.}
\label{fig:bf-covariance-proposal-covariance-roles}
\end{figure}

\section{Conditional local flows and a single proposal mixture}

\subsection{A concrete local affine construction}

A local map should describe uncertainty around one ancestor. The following
construction makes that requirement explicit and supplies a mathematical
realization of the local-flow component. The optional implementation described
at the end of this chapter uses this construction with additive Gaussian
transitions; it does not replace the established LEDH implementation.

For ancestor \(j\), approximate its conditional law by
\(\bfCovNormal(m_j,C_j)\), with \(C_j\in\bfCovSPD^d\).
At a deterministic or independently generated pilot \(\xi_j\), approximate
the log likelihood by a quadratic potential
\begin{equation}
\label{eq:bf-covariance-proposal-quadratic-guide}
 \log \widetilde g_j(x)
 =\mathrm{constant}+a_j^{\rm loc\,\top}x-\tfrac12x^\top B_jx,
 \qquad B_j=B_j^\top\succeq0 .
\end{equation}
Here \(a_j^{\rm loc}\) is a coefficient vector, not the ancestor \(a_j\).
For an additive Gaussian observation, a linearization
\(h(x)\simeq H_jx+e_j\) gives
\(B_j=H_j^\top R^{-1}H_j\) and
\(a_j^{\rm loc}=H_j^\top R^{-1}(y-e_j)\).
The approximation is used to construct a proposal, while \(g_\theta\)
remains the density in the weight.

For a more general differentiable likelihood, one possible score-matching
choice is
\[
 a_j^{\rm loc}
 =\nabla_x\log g_\theta(y\mid \xi_j)+B_j\xi_j .
\]
It matches the local gradient at \(\xi_j\). A specified positive
semidefinite information matrix can supply \(B_j\); for example, a fixed
positive-weight quadrature of observation-score outer products has that
property. The quadrature, simulated observations, weights and their
parameter dependence must be declared. This is a candidate approximation,
not an inherited universal curvature default. Negative-Hessian projection,
if used instead, introduces spectral and differentiability choices that
need separate examination.

A concrete positive-weight construction is
\(B_j=\sum_{r=1}^{R_g}\omega_r s_{jr}s_{jr}^\top\),
where \(\omega_r>0\), \(\sum_r\omega_r=1\), and
\(s_{jr}=\nabla_x\log g_\theta(\widetilde y_{jr}\mid x)|_{x=\xi_j}\).
The pseudo-observations \(\widetilde y_{jr}\) are generated from fixed,
independent base noises by the observation simulator at \(\xi_j\);
their total derivatives are retained. For every vector \(v\),
\(v^\top B_jv=\sum_r\omega_r(v^\top s_{jr})^2\geq0\).
Thus this is a completely defined information approximation when the
model has that simulator and the required derivatives. It is still a
tunable candidate and can be a poor curvature approximation.

Multiplying the Gaussian conditional guide by
\(\widetilde g_j^\lambda\), \(0\leq\lambda\leq1\), and completing the
square gives the following formulas. Both \(C_j\) and
\(C_j^{-1}+\lambda B_j\) are invertible: the first is positive definite
by assumption, and the second is positive definite for
\(0\leq\lambda\leq1\) because \(B_j\succeq0\).
\begin{equation}
\label{eq:bf-covariance-proposal-bridge-moments}
 P_j(\lambda)=(C_j^{-1}+\lambda B_j)^{-1},\qquad
 m_j(\lambda)=P_j(\lambda)
       (C_j^{-1}m_j+\lambda a_j^{\rm loc}).
\end{equation}
The precision is positive definite because \(C_j^{-1}\succ0\) and
\(\lambda B_j\succeq0\). Differentiate with the coefficients held fixed
in pseudo-time:
\begin{equation}
\label{eq:bf-covariance-proposal-bridge-derivatives}
 \dot P_j=-P_jB_jP_j,\qquad
 \dot m_j=P_j(a_j^{\rm loc}-B_jm_j(\lambda)).
\end{equation}
The first identity follows by differentiating
\(P_j(C_j^{-1}+\lambda B_j)=I\). Substituting it into the derivative
of \(m_j(\lambda)\) proves the second.

Define the affine ODE coefficients
\begin{equation}
\label{eq:bf-covariance-proposal-local-ode}
 F_j(\lambda)=-\tfrac12P_j(\lambda)B_j,\qquad
 b_j(\lambda)=\dot m_j(\lambda)-F_j(\lambda)m_j(\lambda),
 \qquad \dot x=F_j(\lambda)x+b_j(\lambda).
\end{equation}
For a Gaussian input, this ODE has the prescribed mean and covariance:
its mean derivative is \(F_jm_j+b_j=\dot m_j\), and its covariance
derivative is
\[
 F_jP_j+P_jF_j^\top
 =-\tfrac12P_jB_jP_j-\tfrac12P_jB_jP_j
 =\dot P_j .
\]
The Gaussian law is determined by these moments. The exact ODE therefore
transports the Gaussian guide to the quadratic-guide posterior.

A fixed discretization defines a concrete invertible proposal map even
when it only approximates this ODE. On a frozen pseudo-time grid, an Euler
step has \(M_k=I+h_kF_j(\lambda_k)\), \(d_k=h_kb_j(\lambda_k)\).
Compose
\begin{equation}
\label{eq:bf-covariance-proposal-affine-composition}
 A^{(0)}=I,\quad c^{(0)}=0,\qquad
 A^{(k+1)}=M_kA^{(k)},\quad
 c^{(k+1)}=M_kc^{(k)}+d_k .
\end{equation}
Every \(M_k\) must be nonsingular with adequate numerical margin.
The final \(T_{L,j}(u)=A_{L,j}u+c_{L,j}\) is then an affine bijection.
Its actual determinant is used in the density; ODE integration error
does not justify replacing it by the determinant of an ideal limiting map.
For fixed \(B_j\), a direct Cholesky map between the endpoint Gaussian
moments is another valid affine realization.

A moving pilot may refresh the local quadratic guide on the grid, provided
the entire pilot design is fixed independently of the particle's transition
noise. The refreshed map remains a proposal with an evaluable density.
The exact Gaussian-bridge statement above applies to fixed coefficients,
not automatically to a sequence of refreshed quadratic approximations.

\subsection{Why retain a likelihood-aware component?}

An ordinary UKF gain uses state--observation cross covariance. It can miss
information conveyed purely through conditional variance. For example, in
the scalar stochastic-volatility observation
\(Y=\exp(X/2)\epsilon\), where \(\epsilon\sim\bfCovNormal(0,1)\) is independent
of \(X\), suitable finite moments give
\(\E[Y\mid X]=0\) and hence \(\Cov(X,Y)=0\).
Nevertheless,
\begin{equation}
\label{eq:bf-covariance-proposal-sv-score}
 \frac{\partial}{\partial x}\log g(y\mid x)
 =\tfrac12(y^2e^{-x}-1)
\end{equation}
is usually nonzero. It follows by differentiating
\(\log g(y\mid x)=-\tfrac12\log(2\pi)-x/2-y^2e^{-x}/2\).
A local proposal informed by this likelihood can respond when the
ordinary UKF mean update cannot. This calculation motivates retaining a
likelihood-aware branch; it does not certify the current KSC-SV call chain
or choose its tuning settings.

\subsection{The actual marginal density}

Use three branch labels \(b\in\{0,G,L\}\):
\(T_{0,j}(u)=u\), \(T_{G,j}=T_G\), and \(T_{L,j}\) as above.
Let \(\beta_b\geq0\), \(\sum_b\beta_b=1\), and \(\beta_0>0\).
These probabilities are frozen within a parameter-evaluation scope.
Section~\ref{sec:bf-covariance-proposal-beta-selection} specifies how to
select them before that scope begins. They are probabilities of choosing
a map, rather than coefficients for averaging the maps themselves.
Conditional on the incoming cloud and any independent pilot, every map
must be fixed before its input transition draw is generated.

\begin{proposition}[Density of the declared sampler]
Sample \(J\) uniformly, \(B\) with probabilities \(\beta_b\), and
\(U\mid J=j\sim f_j\), independently of \(B\) conditional on the fixed
proposal construction. Set \(X=T_{B,J}(U)\). If every map is an invertible
affine map \(T_{b,j}(u)=A_{b,j}u+c_{b,j}\), then
\begin{equation}
\label{eq:bf-covariance-proposal-mapped-density}
 q_{b,j}(x)=
 f_j\bigl(A_{b,j}^{-1}(x-c_{b,j})\bigr)
 |\det A_{b,j}|^{-1},\qquad
 q(x)=\frac1N\sum_j\sum_b\beta_b q_{b,j}(x).
\end{equation}
\end{proposition}
\begin{proof}
The inverse change of variables has Jacobian \(A_{b,j}^{-1}\), proving
the first formula. Conditioning on the sampled pair \((B,J)\) and summing
its probabilities proves the second.
\end{proof}

All branches use one density evaluator with different affine parameters.
They do not require different importance-correction implementations.
If a map were fitted to its own random input \(U\), its derivative would
generally include that fitting dependence. Treating the fitted matrix
\(A(U)\) alone as the Jacobian would then be wrong. Independent pilots
avoid that particular error.

Uniform ancestor selection is appropriate because the incoming reset
weights are equal. If a different route carries nonuniform ancestor
weights, those weights must appear in the sampler, predictive density,
proposal density and their derivatives. Fixed branch allocations can
replace random branch labels, but the denominator must then use the
actual allocation fractions for the pooled estimator.

\section{Importance correction, likelihood and its limits}

For samples \(x_i\) from the declared proposal, define
\begin{equation}
\label{eq:bf-covariance-proposal-importance-ratio}
 r_i=g_\theta(y\mid x_i)\frac{p_N^-(x_i)}{q(x_i)},\qquad
 \widehat Z_N=\frac1N\sum_i r_i,\qquad
 w_i=\frac{r_i}{\sum_k r_k}.
\end{equation}
For logarithmic derivatives we work on a region where all participating
densities and \(\widehat Z_N\) are positive and finite.
Zero target weights require an explicitly supported limiting or masked
calculation, rather than taking an undefined logarithm.

\begin{proposition}[Conditional integration identity and defensive bound]
Under the density and support conditions above,
\begin{equation}
\label{eq:bf-covariance-proposal-is-identity}
 \E_q[r(X)]
 =\int q(x)\frac{g_\theta(y\mid x)p_N^-(x)}{q(x)}\,\bfCovDiff x
 =Z_N .
\end{equation}
Furthermore, wherever the target is positive,
\begin{equation}
\label{eq:bf-covariance-proposal-defensive-bound}
 q(x)\geq\beta_0p_N^-(x),\qquad
 0\leq r(x)\leq\frac{g_\theta(y\mid x)}{\beta_0}.
\end{equation}
\end{proposition}
\begin{proof}
The identity component has density \(q_{0,j}=f_j\). Its contribution to
\(q\) is \(\beta_0p_N^-\); the remaining components are nonnegative.
This proves the lower bound and ensures proposal support wherever the
predictive density is positive. Cancellation inside the integral proves
\eqref{eq:bf-covariance-proposal-is-identity}. Dividing the nonnegative target by the lower
bound proves \eqref{eq:bf-covariance-proposal-defensive-bound}.
\end{proof}

If \(g_\theta(y\mid x)\) is bounded in \(x\), the importance ratio is
bounded. There is no implied bound on its parameter derivative, no
guarantee of high effective sample size, and no guarantee that a finite
draw visits all relevant modes. A poor UKF guide changes the variance of
the integration scheme, not the target in \eqref{eq:bf-covariance-proposal-is-identity}.

The observation increment is
\begin{equation}
\label{eq:bf-covariance-proposal-log-increment}
 \widehat\ell_t
 =\log\widehat Z_{N,t}
 =\bfCovLSE_i(\log r_{t,i})-\log N,\qquad
 \widehat\ell=\sum_{t=1}^T\widehat\ell_t.
\end{equation}
The UKF Gaussian observation normalizer is not multiplied into this
quantity. The observation has already informed the proposal; its actual
density appears once in the importance numerator.

The one-step identity does not establish an unbiased full likelihood.
A deterministic reset changes the finite approximation handed to the
next observation, and logarithms introduce their own bias. In particular,
\(\nabla_\theta\widehat\ell\), the derivative of this finite computation,
and \(\nabla_\theta\log p_\theta(y_{1:T})\), the true model score, are
different targets. Their numerical agreement is a research question.


\section{Selecting the three proposal probabilities}
\label{sec:bf-covariance-proposal-beta-selection}

The three transformations supply three alternative ways to place a particle.
Their probabilities determine how often each alternative is used. Writing
\(q_b=N^{-1}\sum_jq_{b,j}\), the sampling operation is
\begin{equation}
\label{eq:bf-covariance-proposal-beta-sampler}
 B\sim\operatorname{Categorical}(\beta_0,\beta_G,\beta_L),\qquad
 X=T_{B,J}(U),\qquad
 q_\beta(x)=\beta_0q_0(x)+\beta_Gq_G(x)+\beta_Lq_L(x).
\end{equation}
Here \(J\) and \(U\) are sampled as in the density proposition. For example,
probabilities \((0.1,0.6,0.3)\) mean that each draw chooses the identity,
global UKF, or conditional-local map with those probabilities. These numbers
illustrate the sampling rule; they are not calibrated settings. Applying
\(\beta_0T_0+\beta_GT_G+\beta_LT_L\) to each input would instead produce
a different proposal, whose density is not the mixture above.

\subsection{An objective that needs no oracle likelihood}

First hold the parameter, observation, incoming cloud, and all three maps
fixed. Define the unnormalized one-step posterior density
\(a(x)=g_\theta(y\mid x)p_N^-(x)\) and its integral \(Z_N=\int a(x)\,\bfCovDiff x\).
Although \(Z_N\) is unknown, \(a(x)\) and every \(q_b(x)\) can be evaluated.
Assume the relevant second moment is finite. For \(N\) independent draws
from \(q_\beta\),
\begin{equation}
\label{eq:bf-covariance-proposal-beta-second-moment}
 J(\beta)=\int\frac{a(x)^2}{q_\beta(x)}\,\bfCovDiff x,
 \qquad
 \Var(\widehat Z_N\mid\hbox{fixed cloud and maps})
 =\frac{J(\beta)-Z_N^2}{N}.
\end{equation}
Indeed, \(\E_{q_\beta}[(a/q_\beta)^2]=J(\beta)\) and
\(\E_{q_\beta}[a/q_\beta]=Z_N\); independence gives the displayed variance
of the sample mean. Minimizing \(J\) therefore minimizes this conditional
variance without knowing \(Z_N\). Its integrand is a nonnegative constant
divided by a positive affine function of \(\beta\), so \(J\) is convex on
its finite domain. This is the mixture-variance criterion of
\citet[Theorem~3, Section~4, equation~(15), and proof in
Section~7.5]{HeOwen2014mixture}, specialized here to three components and no
control variates. Their mixture probabilities are denoted by \(\alpha\);
their control-variate coefficients denoted by \(\beta\) are absent here.

The independence condition matters. Stratified branch allocations and
dependent SQMC points have a different variance expression, including their
covariances. For those designs, \(J\) remains a proposal-quality criterion;
the displayed equality is not their sampling-variance formula. Likewise,
it concerns a one-step normalizer, rather than the variance or error of the
recursive log likelihood or its score.

\subsection{A deliberate amount of transition protection}

The identity branch has \(q_0=p_N^-\). Choose a protection multiplier
\(B_{\rm safe}>1\) and set \(\epsilon=1/B_{\rm safe}\). Optimize over
\begin{equation}
\label{eq:bf-covariance-proposal-beta-feasible}
 \mathcal B_\epsilon
 =\{\beta:\ \beta_0\geq\epsilon,\ \beta_G,\beta_L\geq0,\
                     \ \beta_0+\beta_G+\beta_L=1\},\qquad
 \frac{p_N^-(x)}{q_\beta(x)}\leq B_{\rm safe}.
\end{equation}
The inequality follows from \(q_\beta\geq\epsilon p_N^-\), wherever the
predictive density is positive. It bounds the density-ratio factor in the
weight by a chosen amount: \(r(x)\leq B_{\rm safe}g_\theta(y\mid x)\).
A smaller multiplier reserves more draws for the transition proposal; a
larger one lets the guided branches receive more draws but weakens that
bound. This is a safety/efficiency tradeoff with a precise meaning, not a
universal numerical constant or a bound on score error. If
\(\int g_\theta(y\mid x)^2p_N^-(x)\,\bfCovDiff x<\infty\), it also makes
\(J(\beta)\) finite throughout \(\mathcal B_\epsilon\).

The multiplier is declared before fitting \(\beta\). A research scope must
justify it by its required ratio protection and a separate robustness
assessment. If several multipliers are investigated, that is an additional
outer calibration exercise with fresh validation; minimizing the same fitting
objective is not a justification for weakening the protection. The other
two probabilities may be zero. Retaining a transition floor does not require
spending particles on a redundant guided branch.

\subsection{Estimate the objective on an independent pilot}

Use a fixed pilot density \(h=(q_0+q_G+q_L)/3\). Generate \(M\) independent
draws \(x_m\) from \(h\), and evaluate all three branch densities and
\(a\) at each draw. The equal pilot allocation samples every branch in
expectation; it does not guarantee representation in a small finite pilot
and is not the selected allocation for the filter. Do not replace \(h\)
by the candidate \(q_\beta\) while continuing to reuse these pilot draws.

For each pilot point define \(d_{mb}=q_b(x_m)/h(x_m)\) and
\(w_m=a(x_m)/h(x_m)\). Then
\begin{equation}
\label{eq:bf-covariance-proposal-beta-pilot}
 \widehat J(\beta)
 =\frac1M\sum_{m=1}^M
       \frac{a(x_m)^2}{h(x_m)q_\beta(x_m)}
 =\frac1M\sum_{m=1}^M\frac{w_m^2}{d_m^\top\beta}.
\end{equation}
A pilot row with zero target density contributes zero, including when its
candidate density is zero. For any fixed \(\beta\), expectation under
\(h\) cancels its denominator
and gives \(\E_h[\widehat J(\beta)]=J(\beta)\). This pointwise statement
does not make the minimum found on the pilot an unbiased estimate of its
out-of-sample performance. Use new draws to assess the selected mixture.
The maps must already be fixed before generating these objective-evaluation
draws. If a map itself needs a fitted pilot, use an earlier, separate pilot
for that fitting, as required by the density proposition.

For calibration over a whole filtering scope, assemble \(C\) fixed contexts
before this optimization. A context contains a parameter value, time and
observation, incoming cloud, and fixed maps. Use incoming clouds from a
declared baseline filter on calibration data, and cover the observation
regimes and parameter region in which the configuration will be evaluated.
The baseline and coverage design must be recorded. They need not be oracle
posterior clouds. Fit a single \(\beta\) for the scope by minimizing
\begin{equation}
\label{eq:bf-covariance-proposal-beta-scope-objective}
 F(\beta)=\sum_{c=1}^C\nu_c\,
       \frac{\widehat J_c(\beta)}{s_c^2},\qquad
 \nu_c>0,\quad\sum_c\nu_c=1,\qquad
 s_c=\frac1{M_c}\sum_m w_{cm}>0.
\end{equation}
Here \(\nu_c\) is a predeclared weight for the represented situation;
equal weights are appropriate only if that is the intended coverage.
The fixed scale \(s_c\) prevents contexts with large normalizers from
dominating merely because of their scale. With one context it does not
affect the minimizer. With several contexts it defines an estimated
relative-second-moment criterion. Because \(s_c\) is random, this ratio is
not an unbiased estimate of \(J_c/Z_{N,c}^2\). It remains convex in \(\beta\)
for the fixed pilot. A zero, non-finite, or numerically unreliable scale
makes calibration unusable; silently discarding a difficult context would
change the calibration question.

\subsection{Solve the small constrained problem and freeze its answer}

Stack the pilot rows across contexts, and write
\(\lambda_{cm}=\nu_cw_{cm}^2/(M_cs_c^2)\). These numbers and the three
density ratios per row stay fixed while choosing \(\beta\). Their derivatives
give an explicit solver:
\begin{equation}
\label{eq:bf-covariance-proposal-beta-derivatives}
 \nabla F=-\sum_{c,m}\frac{\lambda_{cm}d_{cm}}
                                {(d_{cm}^\top\beta)^2},\qquad
 \nabla^2F=2\sum_{c,m}\frac{\lambda_{cm}d_{cm}d_{cm}^\top}
                                {(d_{cm}^\top\beta)^3}\succeq0.
\end{equation}
For any vector \(v\), the Hessian quadratic form is
\[
 v^\top\nabla^2F\,v
 =2\sum_{c,m}\frac{\lambda_{cm}(v^\top d_{cm})^2}
                         {(d_{cm}^\top\beta)^3}\geq0.
\]
Rows with zero target density contribute zero and are omitted from these
sums; all retained rows require positive denominators. Identical branch
densities can make the minimizer nonunique without changing the proposal.

A projected-gradient method suffices for the three probabilities. Let
\(l=(\epsilon,0,0)^\top\). Euclidean projection onto the feasible triangle
has the explicit form
\begin{equation}
\label{eq:bf-covariance-proposal-beta-projection}
 [\Pi_{\mathcal B_\epsilon}(v)]_b
   =l_b+\max(v_b-l_b-\tau,0),\qquad
 \sum_b\max(v_b-l_b-\tau,0)=1-\epsilon.
\end{equation}
Sorting the three entries \(v_b-l_b\) finds the scalar threshold \(\tau\).
This is the simplex projection after subtracting the compulsory transition
mass. At an iterate \(\beta\), let \(g=\nabla F(\beta)\) and
\(b_*\in\arg\min_b g_b\). A stopping certificate is
\begin{equation}
\label{eq:bf-covariance-proposal-beta-gap}
 v_*=l+(1-\epsilon)e_{b_*},\qquad
 \mathcal G(\beta)=g^\top(\beta-v_*),\qquad
 0\leq F(\beta)-\min_{z\in\mathcal B_\epsilon}F(z)
             \leq\mathcal G(\beta).
\end{equation}
The vector \(v_*\) minimizes \(g^\top z\) over the triangle. The last
inequality follows from the supporting-hyperplane inequality for a
differentiable convex function. It certifies optimization of the finite
pilot objective, not closeness to the unknown population optimum. This
projected solver is a specified local recipe; the source paper uses a
different numerical optimizer for its broader problem.

\begin{lstlisting}
Algorithm 0. Offline selection of beta for one declared scope

Inputs: fixed calibration contexts and maps; independent pilot designs;
        B_safe > 1, context weights nu; pilot counts M;
        positive tolerance tol and initial step eta0;
        finite iteration and backtracking budgets

epsilon = 1 / B_safe
for each context c:
    sample x[1:M[c]] from h = (q0 + qG + qL) / 3
    evaluate a, h, and all q_b at every sample
    form d, w, s, and lambda in the displayed formulas
    require support, finite evaluations, and reliable positive s

beta = (epsilon,0,0) + (1-epsilon) * (1,1,1) / 3
repeat within the declared iteration budget:
    evaluate F, g, and gap G(beta)
    if G(beta) <= tol * max(1,F): return beta and fit diagnostics
    eta = eta0
    repeat within the backtracking budget:
        candidate = project(beta - eta*g, B_epsilon)
        delta = candidate - beta
        if F(candidate) <= F(beta) + dot(g,delta)
                           + norm(delta)^2/(2*eta): accept this step
        otherwise eta = eta / 2
    if no step accepted: return calibration failure
    beta = candidate
if no convergence: return calibration failure and last gap

Assess the returned beta on independent validation contexts/draws.
Freeze an accepted beta before the separate untouched claim run.
Algorithm 1 reads this fixed beta; it never reruns this optimizer.
\end{lstlisting}

Evaluate mixtures and squared-weight sums in the log domain. In particular,
compute \(w_{cm}/s_c\) through normalized log weights when forming
\(\lambda_{cm}\); do not exponentiate enormous unnormalized weights first.
Record the protection multiplier, context construction, pilot counts and
seeds, objective, stopping gap, validation results, and selected
probabilities. Pilot sizes and solver tolerances are declared numerical
controls, not constants determined by the probability model.

Calibration and validation sets must be disjoint, and both must be
separate from final claim data. Calibration uses no oracle likelihood or
score. Validation can reject a mixture for numerical failures, poor weight
concentration or unstable replication without knowing an oracle; these
diagnostics alone cannot establish likelihood or score accuracy. When
references are available, check the full recursive value and every score
component in each protected model/horizon scope before admitting the
configuration. A scope change requires its own calibration, rather than
copying a mixture fitted to another model or horizon.

This optimization holds incoming clouds and maps fixed; in a running filter,
changing \(\beta\) changes subsequent clouds. It therefore nominates a
configuration for full recursive validation and does not solve the recursive
filter's global optimization problem. In particular, this criterion does
not optimize the variance of the analytical score. Within the admitted
parameter-evaluation scope, \(\bfCovDiff\beta=0\): keep the probabilities and
base branch-label uniforms fixed throughout all likelihood/score calls.
The maps and their parameter derivatives still vary as specified below.
Refitting the probabilities at each parameter value would define another
algorithm, requiring their derivatives and a treatment of changing labels.
The implementation described below executes Algorithm~0 with independent
pilots and records its optimization gap. Its validation objective measures
proposal-weight variability; downstream likelihood and score accuracy require
the separate recursive comparisons described below.


\section{What the resampling transform should preserve}

After importance weighting, the empirical posterior consists of
\(x_1,\ldots,x_N\) with weights \(w_i>0\), \(\sum_iw_i=1\). Put these states
in the columns of \(X\in\R^{d\times N}\), and let \(W=\diag(w)\). Its mean
and covariance are
\begin{equation}
\label{eq:bf-covariance-proposal-weighted-moments}
 \mu_w=Xw,\qquad
 P_w=\sum_iw_i(x_i-\mu_w)(x_i-\mu_w)^\top
     =X(W-ww^\top)X^\top .
\end{equation}
This is a covariance of a probability measure, with no sample-unbiasedness
factor. An equally weighted replacement therefore uses \(1/N\), not
\(1/(N-1)\).

The reset should preserve these moments if its purpose is to represent
the importance-weighted posterior. The UKF covariance \(P_U^+\) has already
served as a proposal guide. Replacing \(P_w\) by \(P_U^+\) here would impose
a different moment approximation after importance correction. Such a
choice could be investigated, but it is a different algorithm.

\subsection{Why a barycentric transform loses spread}

Let a nonnegative coupling \(\Pi\in\R^{N\times N}\) have row sums \(w\)
and column sums \(\bfCovOne/N\). Define a pair of indices by
\(\Pr(I=i,J=j)=\Pi_{ij}\). The source state \(x_I\) has the weighted
empirical law, while the barycentric output
\(y_j=N\sum_i\Pi_{ij}x_i\) is \(\E[x_I\mid J=j]\). Total covariance gives
\begin{equation}
\label{eq:bf-covariance-proposal-barycentric-loss}
 P_w
 =\Cov\!\left(\E[x_I\mid J]\right)
  +\E\!\left[\Cov(x_I\mid J)\right].
\end{equation}
The second term is positive semidefinite. The covariance of the equal-weight
barycentres is therefore no larger than \(P_w\) in the matrix ordering.
This is the precise spread loss that a second-order transform corrects.

\citet[Definition 3.1 and Sections 4--6]{acevedodewiljesreich2017secondorder} give an ensemble-space
formulation. If \(Z=X\widehat D\), it is sufficient to impose
\begin{equation}
\label{eq:bf-covariance-proposal-transform-constraints}
 \widehat D\bfCovOne=Nw,\quad
 \widehat D^\top\bfCovOne=\bfCovOne,\quad
 (\widehat D-w\bfCovOne^\top)(\widehat D-w\bfCovOne^\top)^\top
       =N(W-ww^\top).
\end{equation}
Indeed, \(Z\bfCovOne/N=Xw=\mu_w\). Its centred matrix is
\(X(\widehat D-w\bfCovOne^\top)\), so the final identity gives covariance \(P_w\).
The second constraint makes each output an affine combination of the inputs;
it is useful structurally, although the first and third identities already
establish these two moments.

Starting from a first-order transform \(D\), set \(B=D-w\bfCovOne^\top\) and
\(\widehat D=D+\Delta\), where \(\Delta=\Delta^\top\) and
\(\Delta\bfCovOne=0\). Expanding the last constraint gives the matrix equation
\begin{equation}
\label{eq:bf-covariance-proposal-riccati-correction}
 B\Delta+\Delta B^\top+\Delta^2
   =N(W-ww^\top)-BB^\top.
\end{equation}
The cross terms follow from \((B+\Delta)(B+\Delta)^\top\).
Writing this equation alone does not specify a unique solver, a root,
or its parameter derivative. We therefore give a complete alternative
realization below, whose moment identities can be checked directly.

\subsection{A complete finite-iteration reset}

Assume \(N>d\) and the two covariance matrices factored below are positive
definite. The particle-count condition is necessary for full rank, but does
not guarantee it. Choose a fixed symmetric positive definite cost metric
\(M_c\), entropic scale \(\varepsilon>0\), and number \(L\) of Sinkhorn
iterations. These are declared numerical choices, not universal defaults.
For the source and reference locations both equal to the current \(x_i\), set
\(c_{ij}=(x_i-x_j)^\top M_c(x_i-x_j)\).
Initialize \(\gamma_j^{(0)}=0\), and perform exactly \(L\) iterations:
\begin{align}
\label{eq:bf-covariance-proposal-sinkhorn-iterations}
 \alpha_i^{(\ell)}
  &=\log w_i-\bfCovLSE_j\!\left(-c_{ij}/\varepsilon+
                                 \gamma_j^{(\ell-1)}\right),\\
 \gamma_j^{(\ell)}
  &=-\log N-\bfCovLSE_i\!\left(-c_{ij}/\varepsilon+
                                 \alpha_i^{(\ell)}\right).
 \nonumber
\end{align}
Here \(\bfCovLSE_j a_j=\log\sum_j e^{a_j}\). Define
\(\log\Pi_{ij}=\alpha_i^{(L)}-c_{ij}/\varepsilon+\gamma_j^{(L)}\).
In exact arithmetic the last update fixes the column sums; the row sums
need not yet equal \(w\). Finite iteration and finite precision are reasons
to compute the actual column masses rather than silently assume both
constraints have converged:
\begin{equation}
\label{eq:bf-covariance-proposal-actual-barycentres}
 d_{ij}=\frac{\Pi_{ij}}{\sum_k\Pi_{kj}},\quad
 y_j=\sum_i d_{ij}x_i,\quad
 \bar y=\frac1N\sum_j y_j,\quad
 S_y=\frac1N\sum_j(y_j-\bar y)(y_j-\bar y)^\top .
\end{equation}
Column softmax evaluation of \(\log\Pi\) avoids unnecessary underflow.
The observation covariance denoted \(S_{\rm obs}\) earlier is a different
quantity from this barycentre covariance \(S_y\).

Let \(L_wL_w^\top=P_w\) and \(L_yL_y^\top=S_y\) be Cholesky factorizations
with positive diagonal entries. Thus \(L_w\) and \(L_y\) are invertible;
in particular, the inverse of \(L_y\) below is well-defined.
The proposed reset is
\begin{equation}
\label{eq:bf-covariance-proposal-cholesky-reset}
 z_j=\mu_w+L_wL_y^{-1}(y_j-\bar y).
\end{equation}
Triangular solves implement the inverse action.

\begin{proposition}[Mean and covariance of the actual reset]
Under the stated full-rank assumptions, the outputs in
\eqref{eq:bf-covariance-proposal-cholesky-reset} satisfy
\begin{equation}
\label{eq:bf-covariance-proposal-reset-covariance}
 \frac1N\sum_jz_j=\mu_w,\qquad
 \frac1N\sum_j(z_j-\mu_w)(z_j-\mu_w)^\top=P_w .
\end{equation}
These identities do not require the finite Sinkhorn iterate to have exact
row sums.
\end{proposition}
\begin{proof}
The centred \(y_j\)'s sum to zero, which proves the first identity.
Their covariance is the actual matrix \(S_y\). The transformed covariance is
\(L_wL_y^{-1}S_yL_y^{-\top}L_w^\top
 =L_wL_y^{-1}L_yL_y^\top L_y^{-\top}L_w^\top
 =L_wL_w^\top=P_w\).
\end{proof}

This construction specifies the mathematics of a candidate reset; it does
not assert that an existing repository callable executes every operation
written here. In particular, adding a ridge to either matrix changes this
identity unless its effect is explicitly accounted for. The basic algorithm
rejects rank failure and records numerical validity separately from the
model likelihood. A failed factorization does not imply that the model's
likelihood is zero. A low-rank extension or regularized target needs its own
derivation and score.

Covariance preservation also does not ensure support preservation.
Consider a weighted scalar law with mass \(0.99\) at zero and mass \(0.01\)
at 100. Its mean is one and its variance is 99. Two equally weighted
nonnegative numbers with mean one have the form \(1-a,1+a\), with
\(|a|\leq1\), so their variance is at most one. No transform can preserve
these moments, the support constraint, and that ensemble size simultaneously.
The negative-coefficient/support issue is also explicit in
\citet[Section 7.3]{acevedodewiljesreich2017secondorder}. A smooth unconstrained chart can be useful
when available, but preserving moments in a chart is not preserving moments
in the original physical coordinates.

\section[Bounded higher-moment correction]{A bounded correction for selected third and fourth moments}

The earlier concern about uncontrolled kurtosis fitting remains relevant.
An unconstrained fit can move a few particles very far even while reducing
a selected moment residual. We propose an optional correction that preserves
the already established mean and covariance algebraically and limits every
particle's displacement. Its feasibility and usefulness still require testing.

Write \(E=Z-\mu_w\bfCovOne^\top\), so \(E\bfCovOne=0\) and \(EE^\top/N=P_w\).
For an orthogonal matrix \(O\in\R^{N\times N}\) that fixes \(\bfCovOne\), define
\begin{equation}
\label{eq:bf-covariance-proposal-orthogonal-repair}
 OO^\top=I,\quad O\bfCovOne=\bfCovOne,\qquad
 Z_{\rm new}=\mu_w\bfCovOne^\top+EO.
\end{equation}
The new mean is \(\mu_w+EO\bfCovOne/N=\mu_w\).
The new covariance is \(EOO^\top E^\top/N=P_w\).
Third and fourth moments are not invariant under this mixing of centred
columns, which is why there is freedom to adjust them.

This is an ensemble-space rotation: it mixes particle columns rather than
rotating physical state coordinates. It moves individual particles while
leaving the physical covariance matrix fixed. For a scalar example, the
four equally weighted points \((-1,-1,1,1)\) have mean zero, variance one,
and standardized fourth moment one. The points
\((-\sqrt2,0,0,\sqrt2)\) have the same mean and variance but fourth moment
two. Both centred row vectors have norm two and lie in the subspace
orthogonal to \(\bfCovOne\); an orthogonal transformation in that subspace
maps one to the other. This demonstrates freedom to change kurtosis,
not feasibility under every displacement cap or convergence of a
particular optimizer. In fact the first ensemble is a stationary
configuration for its fourth moment under smooth variance-preserving
rotations, so a zero-start gradient method can stall there.
Indeed, along a smooth path \(u(t)\) preserving \(\sum_i u_i^2\),
the derivative of \(N^{-1}\sum_i u_i^4\) at this first ensemble is
\(4N^{-1}\sum_i u_i^3\dot u_i
=4N^{-1}\sum_i u_i\dot u_i=0\), because \(u_i^3=u_i\) and
the derivative of the fixed squared norm is zero.

\subsection{A smooth parameterization with explicit displacement bounds}

Choose a skew-symmetric matrix \(K\), \(K^\top=-K\), satisfying \(K\bfCovOne=0\).
The matrix \(H=I-K/2\) is invertible for every such real \(K\), as the
norm identity immediately below proves. Its Cayley transform is
\begin{equation}
\label{eq:bf-covariance-proposal-cayley}
 H=I-K/2,\qquad O=H^{-1}(I+K/2).
\end{equation}
For every real vector \(v\),
\(\|Hv\|^2=\|v\|^2+\|Kv\|^2/4\), since \(v^\top Kv=0\).
Thus \(H\) is nonsingular and \(\|H^{-1}\|_2\leq1\).
The matrices \(I-K/2\) and \(I+K/2\) commute, being polynomials in \(K\).
Together with \(K^\top=-K\), this proves \(OO^\top=I\).
Finally \(H\bfCovOne=\bfCovOne=(I+K/2)\bfCovOne\), so \(O\bfCovOne=\bfCovOne\).

Subtracting the identity yields \(O-I=H^{-1}K\). Consequently,
\begin{equation}
\label{eq:bf-covariance-proposal-cayley-bound}
 \|O-I\|_2\leq\|K\|_2,\qquad
 \|(EO-E)e_i\|_2\leq\|E\|_2\|K\|_2,
\end{equation}
where \(e_i\) selects particle \(i\).
This parameterization excludes orthogonal matrices with eigenvalue \(-1\);
a small displacement cap further restricts what moments can be attained.
There is no assertion that every requested skewness or kurtosis is feasible.

A more interpretable cap uses \(U=L_w^{-1}E\). Here \(L_w\) is invertible
because \(L_wL_w^\top=P_w\succ0\). Since \(UU^\top=NI_d\),
\(\|U\|_2=\sqrt N\). Also the Euclidean norm of row \(k\) of \(E\) is
\(\sqrt{N(P_w)_{kk}}\). Therefore every particle displacement satisfies
\begin{align}
\label{eq:bf-covariance-proposal-dual-caps}
 \|L_w^{-1}(z_i^{\rm new}-z_i)\|_2
       &\leq\sqrt N\,\|K\|_2,\\
 |(z_i^{\rm new}-z_i)_k|
       &\leq\sqrt{N(P_w)_{kk}}\,\|K\|_2 .
 \nonumber
\end{align}
Choose a dimensionless whole-vector radius \(\delta>0\) and dimensionless
coordinate radii \(a_k>0\). Setting
\begin{equation}
\label{eq:bf-covariance-proposal-kappa}
 \kappa=\frac{\min(\delta,a_1,\ldots,a_d)}{\sqrt N}
\end{equation}
and enforcing \(\|K\|_2<\kappa\) ensures both
\(\|L_w^{-1}(z_i^{\rm new}-z_i)\|_2\leq\delta\) and
\(|(z_i^{\rm new}-z_i)_k|\leq a_k\sqrt{(P_w)_{kk}}\).
The bound is conservative but explicit. If physical absolute caps \(c_k\)
are required, replacing \(a_k\) by \(c_k/\sqrt{(P_w)_{kk}}\) introduces
parameter dependence and possible nondifferentiability of the minimum.
That is a different derivative specification from the frozen standardized
radii used here.

One smooth way to enforce the bound is
\begin{equation}
\label{eq:bf-covariance-proposal-bounded-skew}
 J=I-\bfCovOne\bfCovOne^\top/N,\quad
 G(v)=J\{V(v)-V(v)^\top\}J,\quad
 K(v)=\kappa\,\frac{G(v)}{\sqrt{1+\|G(v)\|_F^2}} .
\end{equation}
The map \(V\) is a declared smooth parameterization, for example a linear
combination of fixed basis matrices. Since \(J\bfCovOne=0\), \(G\bfCovOne=0\);
it is also skew-symmetric. Moreover
\(\|K\|_2\leq\|K\|_F<\kappa\).
A low-rank or sparse basis can reduce cost, at the price of less moment
freedom. The choice of basis is part of the candidate, not a hidden
implementation shortcut.

These caps limit the additional rotation step relative to \(Z\).
They do not bound the preceding recolouring move in
\eqref{eq:bf-covariance-proposal-cholesky-reset}. Its conditioning and displacement must be
reported separately. Preventing an explosion in the optional correction
does not make an ill-conditioned earlier reset safe.

\subsection{What is fitted, and what remains fixed}

Since \(P_w\succ0\), its Cholesky factor \(L_w\) is invertible.
Define standardized coordinates \(s(x)=L_w^{-1}(x-\mu_w)\).
For a chosen set \(\mathcal A\) of multiindices with total degree three
or four, let \(s^\nu=\prod_k s_k^{\nu_k}\). The weighted target moments
and an explicit objective are
\begin{equation}
\label{eq:bf-covariance-proposal-moment-objective}
 M_\nu=\sum_iw_i s(x_i)^\nu,\qquad
 F(v)=\frac12\sum_{\nu\in\mathcal A}
 \left\{\frac{N^{-1}\sum_i s(z_i(v))^\nu-M_\nu}{\tau_\nu}\right\}^{\!2},
 \quad \tau_\nu>0 .
\end{equation}
For example, the features may include marginal \(s_k^3,s_k^4\) and
pairwise \(s_k^2s_l,\allowbreak s_ks_l^2,\allowbreak s_k^3s_l,\allowbreak s_ks_l^3,\allowbreak s_k^2s_l^2\).
This specifies what pairwise skew and kurtosis repair means. It does not
match the entire third- or fourth-order tensor unless all required
multiindices are included and the residuals actually vanish.
The scales \(\tau_\nu\) are frozen positive constants chosen during
calibration. The targets \(M_\nu\) are fixed while optimizing \(v\);
they still depend on \(\theta\) when differentiating the full filter.

A complete finite optimizer can use a fixed number \(M_{\rm opt}\) of
gradient steps with fixed step sizes \(\eta_r>0\):
\begin{equation}
\label{eq:bf-covariance-proposal-fixed-optimizer}
 v^{(0)}=0,\qquad
 v^{(r+1)}=v^{(r)}-\eta_r\nabla_v F(v^{(r)}),
 \quad r=0,\ldots,M_{\rm opt}-1.
\end{equation}
Take \(V(0)=0\), so the initial output is the unrotated reset.
Every finite iterate obeys the displacement cap, regardless of whether
the objective decreases. Finite parameter values and all matrix solves
must still pass numerical checks. The algorithm records initial and
final residuals; it does not promise descent for arbitrary step sizes.
Choosing the best iterate or reverting according to a data-dependent
threshold would define another, generally piecewise differentiable
finite program. Such decisions cannot be silently omitted from its score.

The weighted targets themselves can have large sampling error when
importance weights concentrate in the tails. Exact fitting of a noisy
fourth moment need not improve the approximation to the true posterior.
This optional stage therefore needs a dedicated non-harm evaluation:
unchanged first and second moments, respected displacement caps, and
no unacceptable deterioration of likelihoods or scores. It remains
unpromoted until those checks are made.

There is also an asymptotic issue. For any Lipschitz function \(\varphi\)
with constant \(L_\varphi\),
\begin{equation}
\label{eq:bf-covariance-proposal-weak-displacement}
 \left|\frac1N\sum_i\varphi(z_i^{\rm new})
       -\frac1N\sum_i\varphi(z_i)\right|
 \leq L_\varphi\max_i\|z_i^{\rm new}-z_i\|.
\end{equation}
If the right-hand displacement tends to zero, the correction inherits
the same weak limit for such tests as the original reset, provided that
limit exists. A fixed positive cap does not establish that condition.
Covariance identities alone therefore prove neither consistency of this
correction nor consistency of the complete filter.


\section{The analytical score of the same finite computation}
\label{sec:bf-covariance-proposal-analytical-score}

For parameter inference, the derivative must follow the computation that
actually produced \(\widehat\ell\). Replacing a finite optimization loop by
an ideal optimum, holding source weights fixed in the reset, or omitting
a map's parameter dependence would give a different derivative.

Fix all base noises, ancestor and branch labels, sigma-point rule constants,
flow grid points, Sinkhorn iteration counts, optimizer iteration counts,
and calibration choices. Consider a parameter neighbourhood in which the
relevant densities are positive, the maps are invertible, the required
covariances are positive definite, and the declared formulas are
differentiable. We use \(\bfCovDiff\) for a total directional derivative in
\(\theta\). Repeating the calculation for parameter directions gives the
gradient. If a support boundary or validity decision is crossed, these
smooth-region formulas alone do not justify differentiability across it.

\subsection{Likelihood and mixture derivatives}

Differentiating the logarithm of the mean importance ratio gives
\begin{equation}
\label{eq:bf-covariance-proposal-total-score}
 \bfCovDiff\widehat\ell
 =\sum_{t=1}^T\sum_{i=1}^Nw_{t,i}\,\bfCovDiff\log r_{t,i},
 \qquad
 \bfCovDiff\log r_i=\bfCovDiff\log g_\theta(y\mid x_i)
             +\bfCovDiff\log p_N^-(x_i)-\bfCovDiff\log q(x_i).
\end{equation}
For example,
\(\bfCovDiff\log\sum_i r_i=(\sum_i r_i)^{-1}\sum_i r_i\,\bfCovDiff\log r_i\).
There is no additional \(\bfCovDiff w_i\) term in this first derivative;
the weights are the coefficients arising from differentiating log-sum-exp.
They do have derivatives in subsequent particle propagation.

The responsibilities at a state \(x\) are
\begin{equation}
\label{eq:bf-covariance-proposal-responsibilities}
 \rho_j(x)=\frac{f_j(x)}{\sum_k f_k(x)},\qquad
 \eta_{b,j}(x)=
 \frac{\beta_bq_{b,j}(x)}
      {\sum_{c,k}\beta_cq_{c,k}(x)}.
\end{equation}
Differentiating the two log-sums yields
\begin{equation}
\label{eq:bf-covariance-proposal-mixture-differentials}
 \bfCovDiff\log p_N^-(x)=\sum_j\rho_j(x)\,\bfCovDiff\log f_j(x),\qquad
 \bfCovDiff\log q(x)=\sum_{b,j}\eta_{b,j}(x)\,\bfCovDiff\log q_{b,j}(x).
\end{equation}
These are total derivatives, including movement of the evaluated state
and of every ancestor. The \(\beta_b\)'s are frozen. Parameter-dependent
mixture probabilities would introduce their derivatives and a separate
sampling-path issue; that is not the algorithm defined here. Terms with
identically zero branch probability are omitted.

For a component, assume \(A\) is invertible throughout the
parameter neighbourhood. With \(u=A^{-1}(x-c)\), differentiate \(Au=x-c\):
\begin{equation}
\label{eq:bf-covariance-proposal-inverse-tangent}
 \bfCovDiff u=A^{-1}(\bfCovDiff x-\bfCovDiff c-(\bfCovDiff A)u),\qquad
 \bfCovDiff\log|\det A|=\tr(A^{-1}\bfCovDiff A).
\end{equation}
The determinant formula holds on a neighbourhood where \(A\) is
nonsingular; its sign is then locally constant. The component uses the same invertible \(A\).
For \(f_\theta(u\mid a)\),
\begin{equation}
\label{eq:bf-covariance-proposal-component-tangent}
 \begin{aligned}
 \bfCovDiff\log q_{b,j}(x)
 &=(\partial_\theta\log f_\theta(u\mid a))\,\bfCovDiff\theta
  +(\nabla_u\log f_\theta(u\mid a))^\top\bfCovDiff u\\
 &\quad +(\nabla_a\log f_\theta(u\mid a))^\top\bfCovDiff a
  -\tr(A^{-1}\bfCovDiff A).
 \end{aligned}
\end{equation}
The observation term is differentiated similarly, with the observation
held fixed. The actual sampled state has tangent
\(\bfCovDiff x=(\bfCovDiff A)u+A\,\bfCovDiff u_{\rm base}+\bfCovDiff c\), where
\(\bfCovDiff u_{\rm base}\) is obtained by differentiating the transition simulator
with its base noise fixed. Equation \eqref{eq:bf-covariance-proposal-inverse-tangent} is used
when evaluating any mixture component at that state; its \(u\) need not
be the generating component's original input.

\subsection{Prediction, Cholesky factors, and the UKF guide}

Set \(v_j=m_j-m^-\). Differentiating total covariance gives
\begin{equation}
\label{eq:bf-covariance-proposal-prediction-tangent}
 \bfCovDiff m^-=\frac1N\sum_j\bfCovDiff m_j,\qquad
 \bfCovDiff P^-=\frac1N\sum_j\left[
 \bfCovDiff C_j+(\bfCovDiff v_j)v_j^\top+v_j(\bfCovDiff v_j)^\top\right],
 \quad \bfCovDiff v_j=\bfCovDiff m_j-\bfCovDiff m^-.
\end{equation}
Conditional-moment approximations, when used, must be differentiated as
the actual approximations being evaluated.

For any \(P=LL^\top\succ0\) with the positive-diagonal Cholesky factor,
put \(B=L^{-1}(\bfCovDiff P)L^{-\top}\). Define \(\Phi(B)\) to keep the strictly
lower triangle, halve the diagonal, and zero the upper triangle. Then
\begin{equation}
\label{eq:bf-covariance-proposal-chol-tangent}
 \bfCovDiff L=L\Phi(B).
\end{equation}
To verify this, let \(C=L^{-1}\bfCovDiff L\), which is lower triangular.
Differentiation of \(P=LL^\top\) gives \(B=C+C^\top\), whose unique
lower-triangular solution is \(C=\Phi(B)\).

Differentiate every sigma point
\(\chi_r=m^-+L_-\zeta_r\) as
\(\bfCovDiff\chi_r=\bfCovDiff m^-+(\bfCovDiff L_-)\zeta_r\), and then the model evaluations
and the sums defining \(\bar y_{\rm obs},S_{\rm obs},C_{xy}\).
The inverse rule follows from differentiating \(MM^{-1}=I\):
\(\bfCovDiff(M^{-1})=-M^{-1}(\bfCovDiff M)M^{-1}\).
For the UKF formulas, \(S_{\rm obs}\succ0\) makes the inverse
well-defined. The rule gives
\begin{align}
\label{eq:bf-covariance-proposal-ukf-tangent}
 \bfCovDiff K_U&=(\bfCovDiff C_{xy})S_{\rm obs}^{-1}
                -K_U(\bfCovDiff S_{\rm obs})S_{\rm obs}^{-1},\\
 \bfCovDiff m_U^+&=\bfCovDiff m^-+(\bfCovDiff K_U)(y-\bar y_{\rm obs})
                   -K_U\,\bfCovDiff\bar y_{\rm obs},\nonumber\\
 \bfCovDiff P_U^+&=\bfCovDiff P^--(\bfCovDiff K_U)S_{\rm obs}K_U^\top
                    -K_U(\bfCovDiff S_{\rm obs})K_U^\top
                    -K_US_{\rm obs}(\bfCovDiff K_U)^\top.\nonumber
\end{align}
The common affine map uses the invertible Cholesky factor \(L_-\)
of \(P^-\succ0\). It then has tangents
\begin{equation}
\label{eq:bf-covariance-proposal-global-tangent}
 \bfCovDiff A_G=(\bfCovDiff L_+-A_G\bfCovDiff L_-)L_-^{-1},\qquad
 \bfCovDiff c_G=\bfCovDiff m_U^+-(\bfCovDiff A_G)m^- -A_G\bfCovDiff m^-.
\end{equation}

The same rules differentiate each local bridge and every actual flow
step. If \(A_{\rm new}=MA\) and \(c_{\rm new}=Mc+d\), then
\[
 \bfCovDiff A_{\rm new}=(\bfCovDiff M)A+M\bfCovDiff A,\qquad
 \bfCovDiff c_{\rm new}=(\bfCovDiff M)c+M\bfCovDiff c+\bfCovDiff d .
\]
The derivatives of \(B_j,a_j^{\rm loc},C_j\) include their dependence on
ancestors, parameters, and any pilot calculations. A likelihood-score-based
local guide may require higher model derivatives when differentiated.
Calling the filter's derivative \emph{analytical} does not remove these
model-level requirements.

\subsection{Source weights and reset derivatives}

Let \(\bar h=\sum_kw_k\,\bfCovDiff\log r_k\). Normalized weights satisfy
\begin{equation}
\label{eq:bf-covariance-proposal-weight-tangent}
 \bfCovDiff w_i=w_i(\bfCovDiff\log r_i-\bar h),\qquad \sum_i\bfCovDiff w_i=0.
\end{equation}
For \(e_i=x_i-\mu_w\), the weighted moments satisfy
\begin{align}
\label{eq:bf-covariance-proposal-weighted-moment-tangent}
 \bfCovDiff\mu_w
   &=\sum_i\{(\bfCovDiff w_i)x_i+w_i\bfCovDiff x_i\},\\
 \bfCovDiff P_w
   &=\sum_i\left[
      (\bfCovDiff w_i)e_ie_i^\top+
      w_i\{(\bfCovDiff e_i)e_i^\top+e_i(\bfCovDiff e_i)^\top\}\right],
 \quad \bfCovDiff e_i=\bfCovDiff x_i-\bfCovDiff\mu_w.\nonumber
\end{align}
The terms involving \(\bfCovDiff w_i\) and the direct source-cloud derivatives
are required even if the transported barycentres have already been
differentiated. Omitting them is a partial derivative, wrong relative to
the total derivative in \eqref{eq:bf-covariance-proposal-total-score}.

Every finite Sinkhorn update is differentiated through its log-sum-exp.
For example, if \(l_j=\log\sum_i e^{a_{ij}}\), then
\(\bfCovDiff l_j=\sum_i\bfCovSoftmax_i(a_{\cdot j})\,\bfCovDiff a_{ij}\).
With the fixed metric used above,
\[
 \bfCovDiff c_{ij}
 =2(x_i-x_j)^\top M_c(\bfCovDiff x_i-\bfCovDiff x_j).
\]
The actual normalized columns satisfy
\[
 \bfCovDiff d_{ij}
 =d_{ij}\left(\bfCovDiff\log\Pi_{ij}
       -\sum_kd_{kj}\bfCovDiff\log\Pi_{kj}\right),\qquad
 \bfCovDiff y_j=\sum_i\{(\bfCovDiff d_{ij})x_i+d_{ij}\bfCovDiff x_i\}.
\]
Differentiate their actual mean and covariance, use
\eqref{eq:bf-covariance-proposal-chol-tangent} for both factors, and differentiate
\eqref{eq:bf-covariance-proposal-cholesky-reset}. Explicitly, put
\(A_R=L_wL_y^{-1}\); then
\[
 \bfCovDiff A_R=(\bfCovDiff L_w-A_R\bfCovDiff L_y)L_y^{-1},\qquad
 \bfCovDiff z_j=\bfCovDiff\mu_w+(\bfCovDiff A_R)(y_j-\bar y)
                         +A_R(\bfCovDiff y_j-\bfCovDiff\bar y).
\]
Differentiating a converged optimal-transport solution implicitly would
be a different computation from differentiating these \(L\) finite
iterations.

\subsection{Optional higher-moment derivatives}

Since \(K^\top=-K\), the preceding norm identity makes
\(H=I-K/2\) invertible. Differentiation of \(HO=I+K/2\) gives
\begin{equation}
\label{eq:bf-covariance-proposal-cayley-tangent}
 \bfCovDiff O=H^{-1}\frac{\bfCovDiff K}{2}(I+O),\qquad
 \bfCovDiff Z_{\rm new}=(\bfCovDiff\mu_w)\bfCovOne^\top+(\bfCovDiff E)O+E\,\bfCovDiff O .
\end{equation}
For \(s=\sqrt{1+\|G\|_F^2}\), the bounded parameterization gives
\[
 \bfCovDiff K=\frac{\bfCovDiff\kappa}{s}G+\frac{\kappa}{s}\bfCovDiff G
       -\frac{\kappa}{s^3}G\langle G,\bfCovDiff G\rangle_F.
\]
Here \(\bfCovDiff\kappa=0\) for the frozen standardized radii and fixed \(N\).
All dependence of the moments, whitening factor, and feature evaluations
on \(\theta\) remains in \(\bfCovDiff G\) through the optimized parameters
and in the objective derivatives.

Writing \(I_{\rm in}\) for the optimizer's input quantities, the tangent
of an actual finite gradient step is
\begin{equation}
\label{eq:bf-covariance-proposal-optimizer-tangent}
 \bfCovDiff v^{(r+1)}
 =\bfCovDiff v^{(r)}
 -\eta_r\left[
   \nabla_v^2F(v^{(r)};I_{\rm in})\,\bfCovDiff v^{(r)}
   +\bfCovDiff_{I_{\rm in}}\nabla_vF(v^{(r)};I_{\rm in})\right].
\end{equation}
The second term includes the source particles, weights, moment targets,
and Cholesky whitening. There is no assumption that the last iterate is
stationary. The formula may be expensive, but dropping these terms would
change the score.

These recurrences define the analytical derivative of the proposed
finite program. An implementation audit must additionally show that every value/score entry point reaches these same operations,
and verify directional derivatives with common-noise finite differences.
Even exact finite-program derivative parity would not establish agreement
with the true model score.


\section{The whole algorithm in execution order}

The following pseudocode specifies the value computation and the associated
total derivative. Its model interface provides the transition simulator and
density, observation density, conditional moments or a declared approximation
to them, and their required derivatives. The initial sampler is treated in
the same way. For a non-additive observation model, \texttt{UKF\_GUIDE}
uses an augmented sigma-point rule for joint observation/state moments.

The fixed configuration contains the branch probabilities selected by
Algorithm~0 in Section~\ref{sec:bf-covariance-proposal-beta-selection}, the sigma-point rule,
local quadratic-guide construction, flow grid, cost metric, entropic scale,
all iteration counts, and optional moment features, basis, scales, radii,
and step sizes. A user can run a conditional-local-only comparator or a
global-guide ablation by fixing the relevant branch probabilities to zero.
The proposal density always uses the probabilities of the sampler actually
run. The full candidate retains \(\beta_0>0\) for the transition branch.
No model-independent numerical choices for these controls are established
by the derivations.

\begin{lstlisting}
Algorithm 1. One fixed-design likelihood and analytical score

Inputs:
  model, theta, observations y[1:T], particle count N > d
  frozen configuration C and fixed base random design U
Outputs:
  ell, its total derivative score, equal-weight particles,
  per-step diagnostics and an explicit numerical-validity status

(a[1:N], da[1:N]) = INITIAL_SAMPLE_AND_TANGENT(theta, U.initial)
ell = 0
score = 0

for t = 1,...,T:

  # Conditional spread and whole predictive spread are different.
  for j = 1,...,N:
    (m[j], Ccond[j], dm[j], dCcond[j]) =
        CONDITIONAL_MOMENTS_AND_TANGENT(theta, a[j], da[j])
  (mminus, Pminus, dmminus, dPminus) =
      TOTAL_COVARIANCE(m, Ccond, dm, dCcond)

  # The UKF filtered covariance is used HERE.
  (mplus, Pplus, dmplus, dPplus) =
      UKF_GUIDE_AND_TANGENT(mminus, Pminus, y[t], theta, C)
  require Pminus and Pplus positive definite
  LGminus = chol(Pminus)
  LGplus  = chol(Pplus)
  AG = LGplus * inverse(LGminus)
  cG = mplus - AG * mminus
  differentiate these actual factors and solves

  # Build every map before generating its integration inputs.
  for j = 1,...,N:
    map[0,j] = (identity, zero)                     # transition
    map[G,j] = (AG, cG)                             # whole cloud
    map[L,j] = CONDITIONAL_AFFINE_FLOW(
        m[j], Ccond[j], theta, y[t], C, U.pilot[t,j])
    differentiate all map coefficients and pilot dependencies
    require every used map nonsingular

  for i = 1,...,N:
    (b,j) = FIXED_LABEL(U.labels[t,i], C.beta, N)
    (u, du) = TRANSITION_SAMPLE_AND_TANGENT(
        theta, a[j], da[j], U.transition[t,i])
    (x[i], dx[i]) = APPLY_MAP_AND_TANGENT(map[b,j], u, du)

    # One evaluator sums ALL ancestors and ALL used branches.
    logp[i] = logsumexp_j(log f_theta(x[i] | a[j])) - log N
    logq[i] = logsumexp_(b,j)(
        log beta[b] + log MAPPED_DENSITY(x[i], map[b,j], a[j]))
        - log N
    logr[i] = log g_theta(y[t] | x[i]) + logp[i] - logq[i]
    dlogr[i] = TOTAL_DERIVATIVE_OF_THE_SAME_THREE_TERMS
    require relevant densities/derivatives finite and supported

  ell   += logsumexp(logr) - log N
  w      = softmax(logr)
  score += sum_i w[i] * dlogr[i]
  dw[i]  = w[i] * (dlogr[i] - sum_k w[k] * dlogr[k])

  (z, dz, reset_diagnostics) =
      RESET_ACTUAL_WEIGHTED_MOMENTS(x, w, dx, dw, C)

  if C.use_bounded_higher_moment_correction:
    (z, dz, moment_diagnostics) =
        BOUNDED_MOMENT_CORRECTION(x, w, dx, dw, z, dz, C)

  check final mean/covariance against the weighted source
  check support, factor conditioning, and all numerical margins
  save diagnostics, including any invalidity
  if invalid: return INVALID_COMPUTATION with diagnostics

  (a, da) = (z, dz)             # all new weights are 1/N

return ell, score, a, diagnostics
\end{lstlisting}

The reference to total covariance in the listing means
\eqref{eq:bf-covariance-proposal-mixture-moments}; the score recurrences are those in
Section~\ref{sec:bf-covariance-proposal-analytical-score}. If no global branch is used in an ablation, its unused
factorizations can be skipped, with that finite program recorded explicitly.
In the full candidate \(P^-\) and \(P_U^+\) determine \(A_G\), and the
conditional \(C_j\)'s determine the local maps. There is no subtraction
of \(Q\) from a UKF filtered covariance to manufacture either input.

The label routine uses the fixed uniform random numbers with the frozen
probabilities. It does not resample a parameter-dependent set of ancestor
weights: the incoming particles are equally weighted after reset.
Pilot randomness is separate from the integration inputs. Fixing randomness
makes the value a reproducible function of \(\theta\); independence is the
sampling assumption used for the conditional one-step importance identity.

\noindent\begin{minipage}{\linewidth}
\begin{lstlisting}
Algorithm 2. Reset to the actual importance-weighted moments

Inputs: x, w, dx, dw, and fixed cost metric, epsilon, L
(mu, Pw, dmu, dPw) = WEIGHTED_MOMENTS_AND_TANGENT(x,w,dx,dw)
require Pw positive definite

cost[i,j] = (x[i]-x[j])' * M_cost * (x[i]-x[j])
differentiate the actual cost
gamma = 0
for l = 1,...,L:
  alpha[i] = log w[i] - logsumexp_j(-cost[i,j]/epsilon + gamma[j])
  gamma[j] = -log N - logsumexp_i(-cost[i,j]/epsilon + alpha[i])
  propagate derivatives through both finite updates

logPi[i,j] = alpha[i] - cost[i,j]/epsilon + gamma[j]
D[:,j] = softmax_i(logPi[:,j])            # actual column masses
y[j] = sum_i D[i,j] * x[i]
ybar = mean_j y[j]
Sy = mean_j (y[j]-ybar) * (y[j]-ybar)'
propagate derivatives of D, y, ybar, and Sy
require Sy positive definite

Lw = chol(Pw)
Ly = chol(Sy)
AR = Lw * inverse(Ly)
z[j] = mu + AR * (y[j]-ybar)
propagate derivatives of both Cholesky factors and this map

record actual row/column marginal residuals, factor margins,
reset displacement, and final moment residuals
return z, dz, diagnostics
\end{lstlisting}
\end{minipage}

\noindent\begin{minipage}{\linewidth}
\begin{lstlisting}
Algorithm 3. Optional bounded pairwise/higher-moment correction

Inputs: weighted source (x,w), reset z, all tangents, fixed controls
(mu, Pw, Lw) = source weighted mean, covariance, Cholesky factor
E = z - mu * ones'
target[nu] = sum_i w[i] * standardized(x[i])^nu
v = 0
dv = 0

for r = 0,...,Mopt-1:
  G = J * (V(v) - V(v)') * J
  K = kappa * G / sqrt(1 + frobenius_norm(G)^2)
  O = solve(identity - K/2, identity + K/2)
  candidate = mu * ones' + E * O
  F = scaled squared moment residual against target
  v_next = v - eta[r] * gradient_v(F)
  dv_next = derivative of this actual update,
            including input and objective dependencies
  (v,dv) = (v_next,dv_next)

recompute G, K, O, final z, and all tangents at the final v
record initial/final feature residuals and optimizer finiteness
verify displacement caps and final first/second moments
return final z, dz, diagnostics
\end{lstlisting}
\end{minipage}

All three routines return invalidity rather than silently clipping a state,
adding a covariance ridge, or changing a step count. Such protections may be
valuable, but they alter the finite computation and require declared
semantics, a derivative treatment, and dedicated evaluation. In particular,
\texttt{INVALID\_COMPUTATION} must not be converted to a mathematical assertion
that the observation or parameter has zero probability.

The mixture evaluation is generally quadratic in particle count: each of
\(N\) states is assessed under all ancestors and proposal branches.
Dense transport is also quadratic. Full ensemble-space rotations can cost
more unless a suitable basis and linear algebra reduce them. These costs
are part of the proposal's tradeoff. The covariance formulas do not remove
the need for an implementation and memory analysis.

\section[Literature and proposed extensions]{What the literature supplies, and what is proposed here}

The unscented particle filter of \citet{vandermerwe2000unscented} uses unscented
Gaussian approximations to construct measurement-informed proposals,
followed by importance correction. Its relevance is the role of the UKF
as a proposal construction tool. It does not make a UKF covariance an
independent posterior truth. The common-map/conditional-map mixture
specified here is a separate synthesis.

The second-order ensemble-transform work of \citet{acevedodewiljesreich2017secondorder} directly
answers the spread-loss problem. Its target is the weighted ensemble's
first two moments. Sections 4--6 develop corrections and their solution;
Section 7.3 makes clear that moment constraints can conflict with bounded
state support. These results motivate preserving \(P_w\), while leaving
the accuracy of the weighted ensemble itself to importance sampling and
subsequent validation.

\citet[Algorithm 3, Section 4 and Proposition 4.3]{corenflos2021differentiable} develop
differentiable particle filtering through entropy-regularized transport.
The paper explains why neglecting resampling derivatives gives a different
gradient, and explicitly discusses bias of the finite likelihood estimate.
Its consistency results impose assumptions on the model, potentials,
transport, and limiting regime. They cannot be transferred automatically
to this mixture proposal, Cholesky reset, and bounded finite optimizer.
The author FilterFlow transport/Sinkhorn source was also inspected in the
preceding source review; that corroborates those operations, not this new
combination.

A different response to missing ensemble spread is the covariance-shrinkage
rejuvenation studied by \citet[Sections 4--6]{popov2022rejuvenation}. It augments an
ensemble with synthetic anomalies carrying external climatological covariance.
This can supply directions that the current finite cloud lacks. It also adds
information and changes the represented ensemble, so it is not a proof that
the current weighted cloud already contains the correct posterior covariance.
The paper was inspected; its linked author-code archive was not successfully
retrieved in this review.

For a more flexible observation update, \citet[Sections 3.3, 3.7 and
Appendix A]{spantini2022coupling} construct nonlinear triangular coupling maps.
The inspected author \texttt{StochasticMapFilter.m} fits a map to simulated
observation/state pairs and applies an inverse conditional transformation.
This is a principled way to move beyond Gaussian conditional approximations.
It introduces map-fitting error and cost; differentiating a fitted map
requires a corresponding specification. It is a longer-term alternative,
not a free covariance correction inside the present algorithm.

\citet[Sections III--V]{csuzdi2024optimalplacement} give a differentiable
optimal-placement resampling construction in one dimension.
Its inspected scope does not supply an 18-dimensional replacement.
\citet[Sections IV--V]{ebeigbe2025genut} construct generalized unscented points
using marginal third and fourth moments, and discuss the effect of bound
constraints. Their weighted sigma-point construction does not prove that
an arbitrary equally weighted ensemble can match every mixed moment.
The author GenUT code was checked for its relevant ensemble construction.

Taken together, these works support distinctions and individual mechanisms:
measurement-informed proposals, correction of resampling spread loss,
differentiation through transport, and explicit treatment of higher moments.
The three-branch common/conditional proposal, the particular finite
Cholesky realization, and the bounded Cayley fitting scheme are the
candidate synthesis in this note. No cited paper proves that the combined
algorithm cures the observed SIR likelihood or score errors.

\section[Implementation and promotion checks]{What implementation and promotion require}

A correct density formula can coexist with a useless proposal. A correct
reset can faithfully preserve inaccurate weighted moments. A correct
finite-program score can disagree strongly with the true score. These
are different failure mechanisms and require different comparisons.

The implementation separates the conditional-covariance proposal, the
common UKF branch, and the optional higher-moment optimizer so that each
can be checked independently. Every value and score entry point must reach
the same marginal-density evaluator. Comparisons with the transition-only,
equal-mixture, global-heavy and local-heavy proposals help identify which
component causes a failure or an observed gain.

For the proposal, exact linear-Gaussian identities, density normalization
checks, and common-noise directional finite differences test mathematical
and implementation mechanics. A scalar stochastic-volatility test checks
the observation-cross-covariance blind spot discussed above. The zero-flow
branch and single-component cases test the shared evaluator. These checks
are necessary but do not establish nonlinear filtering accuracy.

For the reset, check the weighted mean and covariance after the \emph{last}
operation, not just after recolouring. Report actual marginal residuals,
relative factorization margins, smallest singular values of flow maps,
state support, reset displacement, moment-feature residuals, and the two
caps. Exact algebra permits floating-point residuals; tolerances need a
scale-aware numerical analysis rather than a universal absolute constant.
A failed support or factorization check is a numerical-validity veto,
whereas effective sample size or a disagreement between \(P_w\) and
\(P_U^+\) is explanatory unless a reviewed experiment defines a different
role. Neither is an oracle-free certificate of likelihood or score accuracy.

Scientific comparison then needs the same prepared data, parameter points,
horizons, particle counts, and uncertainty accounting across eligible
implementations. Protected scopes include LGSSM, KSC stochastic volatility,
range bearing, predator--prey, and SIR, with the observation horizon part of
each tuning scope. Linear-Gaussian values and scores can use exact references.
Other models need checked independent references, including source-grounded
Zhao--Cui computations where applicable, and convergence/uncertainty analysis
of a large bootstrap filter. Previously invalidated historical results
cannot silently become current baselines.

The comparison should also include simple competitors constructed for
the situation: the transition/bootstrap proposal when observations are
weak, the exact conditional Gaussian proposal when available, and a plain
UKF proposal when the joint distribution is close to Gaussian.
Regimes with informative observations, non-Gaussian observation dependence,
concentrated weights, or support boundaries should be assessed separately.
A pooled average can conceal exactly the pathology being investigated.

Branch probabilities, local curvature rules, transport regularization,
caps, feature scales, and optimization step sizes are hypotheses to
calibrate, not consequences of the covariance identities.
Oracle-free diagnostics can nominate settings on separate calibration data;
they cannot prove proximity to the true likelihood or score.
Settings must be frozen before held-out evaluation and parameter-dependent
retuning inside an HMC trajectory is outside this proposal.
The optional correction additionally needs a non-harm study of its caps
and failure handling, including the symmetric stationary example above.
The bounded implementation study below was authorized separately; it does
not establish a generally tuned configuration.

The established results in this note are conditional mathematical statements:
the mixture covariance decomposition, the change-of-variables proposal
density and one-step importance identity, the reset's first two moments,
the rotation's invariants and displacement bounds, and the derivative
recurrences for the stated finite program. Global accuracy, weak convergence
of the entire composite method, reliable tail moments, efficient computation,
and suitability for HMC remain to be demonstrated.


\section{Implementation and first numerical checks}

The optional TensorFlow implementation in
\texttt{covariance\_proposal\_tf.py} executes Algorithms~1 and~2 through one
shared value and analytical directional-derivative function. The companion
\texttt{covariance\_proposal\_beta\_tf.py} implements Algorithm~0, and
\texttt{covariance\_proposal\_moments\_tf.py} implements Algorithm~3.
The master program is
\path{docs/benchmarks/run_ledh_covariance_proposal.py}.
It supplies calibration, execution, test and reporting commands; its saved
results include the configuration, fixed seeds, source hashes, hardware,
likelihood, every parameter score and numerical diagnostics.

This first implementation covers additive Gaussian transitions with a common
conditional covariance, using the conditional mean as each local linearization
point. The actual observation density enters the importance weights, including
the non-Gaussian KSC mixture. The common guide uses the declared unscented rule.
The dense reset supports \(d<N\leq3000\). The optional moment fit uses four
fixed cosine-pair skew directions, unit feature scales, all marginal third and
fourth powers, and all pairwise powers of total degree three or four. Its common
standardized radius sets both the whole-vector and coordinate caps.
These are explicit choices for a first candidate, rather than a full
implementation of every model interface and parameterization allowed above.

The bounded study of 8 October 2026 used LGSSM, KSC stochastic volatility,
predator--prey and \(d=18\) SIR. Independent pilots at the first, middle and
last observation selected \(\beta\) with transition floor \(0.1\), while the
Euler grid of 16 steps and 40 Sinkhorn iterations remained fixed hypotheses.
This stage therefore tunes the branch probabilities only. It does not complete
the separate calibration of every numerical control. CPU reference tests cover
all parameters, both with and without the optional moment correction; GPU/XLA
runs check the same shared implementation.

The tests distinguish three questions. Finite differences check whether the
analytical recurrence differentiates the implemented likelihood. Mean,
covariance and displacement checks test the reset and correction invariants.
Exact Kalman, refined KSC and independent bootstrap comparisons assess proximity
to the model likelihood and score. Passing the first two does not settle the
third. The bootstrap reference also has Monte Carlo error and finite-particle
bias; it is not an exact oracle.

An initial SIR check exposed a precision problem: FP32 with TF32 failed the
final covariance check, while the same saved inputs in FP32 without TF32 passed.
The broader comparison therefore used the explicit FP64 GPU/XLA reference
configuration. At \(T=50\), \(N=128\), one fitted SIR mixture selected
\((\beta_0,\beta_G,\beta_L)=(0.1,0,0.9)\). Its two log likelihoods were
\(-1679.56\) and \(-1678.66\), versus a bootstrap estimate of \(-1677.30\)
with estimated Monte Carlo standard error \(0.36\). The corresponding third
scores were \(-11.80\) and \(-20.84\), versus \(-20.00\) with standard error
\(1.77\). The equal mixture produced an extreme third score of \(-2193.07\)
on one seed, and the global-heavy proposal failed numerical checks after
three or four observations. These outcomes require investigation of recursive
sensitivity and precision. Two candidate seeds and an imperfect reference do
not establish a ranking or suitability for HMC.

The complete values, all score coordinates, finite-difference diagnostics,
failed attempts and source-to-equation audit are preserved in
\begin{quote}\small
\path{docs/benchmarks/}\\
\path{ledh-covariance-proposal-results-20261008.md}
\end{quote}
MathDevMCP found the requested operations and terms but returned unverified
obligations, including parser limitations. The mathematical derivations,
executable invariants and derivative tests supply partial, explicitly scoped
evidence; they do not constitute a formal certificate for the whole filter.
